% Options for packages loaded elsewhere
\PassOptionsToPackage{unicode}{hyperref}
\PassOptionsToPackage{hyphens}{url}
%
\documentclass[
]{article}
\usepackage{amsmath,amssymb}
\usepackage{iftex}
\ifPDFTeX
  \usepackage[T1]{fontenc}
  \usepackage[utf8]{inputenc}
  \usepackage{textcomp} % provide euro and other symbols
\else % if luatex or xetex
  \usepackage{unicode-math} % this also loads fontspec
  \defaultfontfeatures{Scale=MatchLowercase}
  \defaultfontfeatures[\rmfamily]{Ligatures=TeX,Scale=1}
\fi
\usepackage{lmodern}
\ifPDFTeX\else
  % xetex/luatex font selection
\fi
% Use upquote if available, for straight quotes in verbatim environments
\IfFileExists{upquote.sty}{\usepackage{upquote}}{}
\IfFileExists{microtype.sty}{% use microtype if available
  \usepackage[]{microtype}
  \UseMicrotypeSet[protrusion]{basicmath} % disable protrusion for tt fonts
}{}
\makeatletter
\@ifundefined{KOMAClassName}{% if non-KOMA class
  \IfFileExists{parskip.sty}{%
    \usepackage{parskip}
  }{% else
    \setlength{\parindent}{0pt}
    \setlength{\parskip}{6pt plus 2pt minus 1pt}}
}{% if KOMA class
  \KOMAoptions{parskip=half}}
\makeatother
\usepackage{xcolor}
\usepackage[margin=1in]{geometry}
\usepackage{color}
\usepackage{fancyvrb}
\newcommand{\VerbBar}{|}
\newcommand{\VERB}{\Verb[commandchars=\\\{\}]}
\DefineVerbatimEnvironment{Highlighting}{Verbatim}{commandchars=\\\{\}}
% Add ',fontsize=\small' for more characters per line
\usepackage{framed}
\definecolor{shadecolor}{RGB}{248,248,248}
\newenvironment{Shaded}{\begin{snugshade}}{\end{snugshade}}
\newcommand{\AlertTok}[1]{\textcolor[rgb]{0.94,0.16,0.16}{#1}}
\newcommand{\AnnotationTok}[1]{\textcolor[rgb]{0.56,0.35,0.01}{\textbf{\textit{#1}}}}
\newcommand{\AttributeTok}[1]{\textcolor[rgb]{0.13,0.29,0.53}{#1}}
\newcommand{\BaseNTok}[1]{\textcolor[rgb]{0.00,0.00,0.81}{#1}}
\newcommand{\BuiltInTok}[1]{#1}
\newcommand{\CharTok}[1]{\textcolor[rgb]{0.31,0.60,0.02}{#1}}
\newcommand{\CommentTok}[1]{\textcolor[rgb]{0.56,0.35,0.01}{\textit{#1}}}
\newcommand{\CommentVarTok}[1]{\textcolor[rgb]{0.56,0.35,0.01}{\textbf{\textit{#1}}}}
\newcommand{\ConstantTok}[1]{\textcolor[rgb]{0.56,0.35,0.01}{#1}}
\newcommand{\ControlFlowTok}[1]{\textcolor[rgb]{0.13,0.29,0.53}{\textbf{#1}}}
\newcommand{\DataTypeTok}[1]{\textcolor[rgb]{0.13,0.29,0.53}{#1}}
\newcommand{\DecValTok}[1]{\textcolor[rgb]{0.00,0.00,0.81}{#1}}
\newcommand{\DocumentationTok}[1]{\textcolor[rgb]{0.56,0.35,0.01}{\textbf{\textit{#1}}}}
\newcommand{\ErrorTok}[1]{\textcolor[rgb]{0.64,0.00,0.00}{\textbf{#1}}}
\newcommand{\ExtensionTok}[1]{#1}
\newcommand{\FloatTok}[1]{\textcolor[rgb]{0.00,0.00,0.81}{#1}}
\newcommand{\FunctionTok}[1]{\textcolor[rgb]{0.13,0.29,0.53}{\textbf{#1}}}
\newcommand{\ImportTok}[1]{#1}
\newcommand{\InformationTok}[1]{\textcolor[rgb]{0.56,0.35,0.01}{\textbf{\textit{#1}}}}
\newcommand{\KeywordTok}[1]{\textcolor[rgb]{0.13,0.29,0.53}{\textbf{#1}}}
\newcommand{\NormalTok}[1]{#1}
\newcommand{\OperatorTok}[1]{\textcolor[rgb]{0.81,0.36,0.00}{\textbf{#1}}}
\newcommand{\OtherTok}[1]{\textcolor[rgb]{0.56,0.35,0.01}{#1}}
\newcommand{\PreprocessorTok}[1]{\textcolor[rgb]{0.56,0.35,0.01}{\textit{#1}}}
\newcommand{\RegionMarkerTok}[1]{#1}
\newcommand{\SpecialCharTok}[1]{\textcolor[rgb]{0.81,0.36,0.00}{\textbf{#1}}}
\newcommand{\SpecialStringTok}[1]{\textcolor[rgb]{0.31,0.60,0.02}{#1}}
\newcommand{\StringTok}[1]{\textcolor[rgb]{0.31,0.60,0.02}{#1}}
\newcommand{\VariableTok}[1]{\textcolor[rgb]{0.00,0.00,0.00}{#1}}
\newcommand{\VerbatimStringTok}[1]{\textcolor[rgb]{0.31,0.60,0.02}{#1}}
\newcommand{\WarningTok}[1]{\textcolor[rgb]{0.56,0.35,0.01}{\textbf{\textit{#1}}}}
\usepackage{graphicx}
\makeatletter
\newsavebox\pandoc@box
\newcommand*\pandocbounded[1]{% scales image to fit in text height/width
  \sbox\pandoc@box{#1}%
  \Gscale@div\@tempa{\textheight}{\dimexpr\ht\pandoc@box+\dp\pandoc@box\relax}%
  \Gscale@div\@tempb{\linewidth}{\wd\pandoc@box}%
  \ifdim\@tempb\p@<\@tempa\p@\let\@tempa\@tempb\fi% select the smaller of both
  \ifdim\@tempa\p@<\p@\scalebox{\@tempa}{\usebox\pandoc@box}%
  \else\usebox{\pandoc@box}%
  \fi%
}
% Set default figure placement to htbp
\def\fps@figure{htbp}
\makeatother
\setlength{\emergencystretch}{3em} % prevent overfull lines
\providecommand{\tightlist}{%
  \setlength{\itemsep}{0pt}\setlength{\parskip}{0pt}}
\setcounter{secnumdepth}{-\maxdimen} % remove section numbering
\usepackage{bookmark}
\IfFileExists{xurl.sty}{\usepackage{xurl}}{} % add URL line breaks if available
\urlstyle{same}
\hypersetup{
  pdftitle={Midterm Exam},
  pdfauthor={Pablo E. Gutiérrez-Fonseca},
  hidelinks,
  pdfcreator={LaTeX via pandoc}}

\title{Midterm Exam}
\author{Pablo E. Gutiérrez-Fonseca}
\date{2025-Oct-08, 06:48 PM}

\begin{document}
\maketitle

\subsubsection{General Instructions / Instrucciones
Generales}\label{general-instructions-instrucciones-generales}

\begin{itemize}
\item
  This is the Midterm-Exam. Each question is worth 20 points. Este es el
  examen parcial. Cada pregunta vale 20 puntos.
\item
  You must choose 5 out of the 6 questions to answer. \textbf{Your
  maximum score is therefore 100 points.} Debe escoger 5 de las 6
  preguntas para contestar. Su puntaje máximo será 100 puntos.
\item
  It is highly recommended that all your results, including tables and
  outputs, are always visible in your R Markdown document. Es altamente
  recomendado que todos los resultados, incluidas las tablas y salidas,
  sean siempre visibles en su documento de R Markdown.
\item
  The exam must be completed and submitted during class time, on
  \textbf{Thursday, October 9, 2025, from 3:00 to 6:00 PM.} El examen se
  debe realizar y entregar durante el tiempo de clase, es decir, jueves
  9 de octubre, de 3 a 6 pm.
\item
  Use R and the tidyverse packages to answer all questions. Utilice R y
  los paquetes tidyverse para contestar todas las preguntas.
\item
  Work independently and make sure your code is reproducible. Trabaje de
  manera independiente y asegúrese de que su código sea reproducible.
\end{itemize}

\begin{Shaded}
\begin{Highlighting}[]
\FunctionTok{library}\NormalTok{(tidyverse)}
\FunctionTok{library}\NormalTok{(dplyr)}
\FunctionTok{library}\NormalTok{(aod)}
\end{Highlighting}
\end{Shaded}

\begin{Shaded}
\begin{Highlighting}[]
\NormalTok{msleep }\OtherTok{\textless{}{-}} \FunctionTok{glimpse}\NormalTok{(msleep)}
\NormalTok{storms }\OtherTok{\textless{}{-}}\NormalTok{ storms}
\NormalTok{lizards }\OtherTok{\textless{}{-}}\NormalTok{ lizards}
\end{Highlighting}
\end{Shaded}

\begin{Shaded}
\begin{Highlighting}[]
\FunctionTok{head}\NormalTok{(msleep)}
\end{Highlighting}
\end{Shaded}

\begin{verbatim}
## # A tibble: 6 x 11
##   name    genus vore  order conservation sleep_total sleep_rem sleep_cycle awake
##   <chr>   <chr> <chr> <chr> <chr>              <dbl>     <dbl>       <dbl> <dbl>
## 1 Cheetah Acin~ carni Carn~ lc                  12.1      NA        NA      11.9
## 2 Owl mo~ Aotus omni  Prim~ <NA>                17         1.8      NA       7  
## 3 Mounta~ Aplo~ herbi Rode~ nt                  14.4       2.4      NA       9.6
## 4 Greate~ Blar~ omni  Sori~ lc                  14.9       2.3       0.133   9.1
## 5 Cow     Bos   herbi Arti~ domesticated         4         0.7       0.667  20  
## 6 Three-~ Brad~ herbi Pilo~ <NA>                14.4       2.2       0.767   9.6
## # i 2 more variables: brainwt <dbl>, bodywt <dbl>
\end{verbatim}

\begin{Shaded}
\begin{Highlighting}[]
\FunctionTok{head}\NormalTok{(storms)}
\end{Highlighting}
\end{Shaded}

\begin{verbatim}
## # A tibble: 6 x 13
##   name   year month   day  hour   lat  long status       category  wind pressure
##   <chr> <dbl> <dbl> <int> <dbl> <dbl> <dbl> <fct>           <dbl> <int>    <int>
## 1 Amy    1975     6    27     0  27.5 -79   tropical de~       NA    25     1013
## 2 Amy    1975     6    27     6  28.5 -79   tropical de~       NA    25     1013
## 3 Amy    1975     6    27    12  29.5 -79   tropical de~       NA    25     1013
## 4 Amy    1975     6    27    18  30.5 -79   tropical de~       NA    25     1013
## 5 Amy    1975     6    28     0  31.5 -78.8 tropical de~       NA    25     1012
## 6 Amy    1975     6    28     6  32.4 -78.7 tropical de~       NA    25     1012
## # i 2 more variables: tropicalstorm_force_diameter <int>,
## #   hurricane_force_diameter <int>
\end{verbatim}

\begin{Shaded}
\begin{Highlighting}[]
\FunctionTok{head}\NormalTok{(lizards)}
\end{Highlighting}
\end{Shaded}

\begin{verbatim}
##   Site Diameter Height    Time grahami opalinus
## 1  Sun   D <= 2  H < 5   Early      20        2
## 2  Sun   D <= 2  H < 5 Mid-day       8        1
## 3  Sun   D <= 2  H < 5    Late       4        4
## 4  Sun   D <= 2 H >= 5   Early      13        0
## 5  Sun   D <= 2 H >= 5 Mid-day       8        0
## 6  Sun   D <= 2 H >= 5    Late      12        0
\end{verbatim}

\subsubsection{\texorpdfstring{\textbf{Mammals sleep dataset
\texttt{(msleep)}}}{Mammals sleep dataset (msleep)}}\label{mammals-sleep-dataset-msleep}

\begin{enumerate}
\def\labelenumi{\arabic{enumi}.}
\tightlist
\item
  \textbf{Use \texttt{filter()} to subset msleep to only animals who are
  awake for at least 12 hours of the day.} Usa filter() para seleccionar
  del dataset msleep solo los animales que están despiertos al menos 12
  horas al día.
\end{enumerate}

\begin{itemize}
\tightlist
\item
  \textbf{Hint:} ``At least'' like saying ``greater than or equal to,''
  i.e.~\texttt{\textgreater{}=}
\item
  \textbf{Sugerencia:} ``Al menos'' significa mayor o igual a
  (\textgreater=).
\end{itemize}

\begin{enumerate}
\def\labelenumi{\arabic{enumi}.}
\setcounter{enumi}{1}
\tightlist
\item
  \textbf{Use mutate() to create a new column called
  \texttt{percent\_day\_awake} that gives the percentage of the day that
  each species spends awake. Use \texttt{select()} with
  \texttt{everything()} at the end to verify your calculation.} Usa
  mutate() para crear una nueva columna llamada percent\_day\_awake que
  muestre el porcentaje del día que cada especie permanece despierta.
  Usa select() con everything() al final para comprobar que tus cálculos
  son correctos.
\end{enumerate}

\begin{itemize}
\tightlist
\item
  \textbf{Hint:} A day has 24 hours. The column awake indicates the
  average hours a species is awake per day. Calculate the percentage
  with \texttt{(awake\ /\ 24)\ *\ 100}.
\item
  \textbf{Sugerencia:} Un día tiene 24 horas. La columna awake indica
  cuántas horas al día, en promedio, está despierta cada especie.
  Calcula el porcentaje con (awake / 24) * 100.
\end{itemize}

\subsubsection{\texorpdfstring{\textbf{Storm tracks data
\texttt{(storms)}}}{Storm tracks data (storms)}}\label{storm-tracks-data-storms}

\begin{enumerate}
\def\labelenumi{\arabic{enumi}.}
\setcounter{enumi}{2}
\tightlist
\item
  \textbf{Use \texttt{group\_by()} and \texttt{summarise()} to calculate
  the number of storms per year.} Agrupa los datos por año y calcula el
  número de tormentas que ocurrieron cada año.
\end{enumerate}

\begin{itemize}
\tightlist
\item
  \textbf{Hint} group\_by(year) groups the rows by year, and
  summarise(num\_storms = n()) counts how many rows (storms) there are
  per year.
\item
  \textbf{Sugerencia:} group\_by(year) agrupa las filas por año, y
  summarise(num\_storms = n()) cuenta cuántas filas (tormentas) hay por
  año.
\end{itemize}

\begin{enumerate}
\def\labelenumi{\arabic{enumi}.}
\setcounter{enumi}{3}
\tightlist
\item
  \textbf{Use filter() to subset storms to only storms with wind speeds
  greater than 50 knots, then use arrange() to sort by wind descending.}
  Filtra el dataset para quedarte solo con tormentas que tengan
  velocidad de viento mayor a 50 nudos, y luego ordénalas de mayor a
  menor velocidad de viento.
\end{enumerate}

\begin{itemize}
\tightlist
\item
  \textbf{Hint} Use \texttt{filter(wind\ \textgreater{}\ 50)} to select
  strong storms, and \texttt{arrange(desc(wind))} to sort the results
  from highest to lowest.
\item
  \textbf{Sugerencia:} Usa \texttt{filter(wind\ \textgreater{}\ 50)}
  para seleccionar las tormentas fuertes, y \texttt{arrange(desc(wind))}
  para ordenar los resultados de mayor a menor.
\end{itemize}

\subsubsection{\texorpdfstring{\textbf{A Comparison of Site Preferences
of Two Species of Lizard
\texttt{(lizards)}}}{A Comparison of Site Preferences of Two Species of Lizard (lizards)}}\label{a-comparison-of-site-preferences-of-two-species-of-lizard-lizards}

\begin{enumerate}
\def\labelenumi{\arabic{enumi}.}
\setcounter{enumi}{4}
\tightlist
\item
  \textbf{Determine which \texttt{Site} and \texttt{Diameter}
  combination has the \texttt{highest\ total\ grahami}.} Determina cuál
  combinación de \texttt{Sitio} y \texttt{Diámetro} tiene el
  \texttt{total\ de\ grahami\ más\ alto}.
\end{enumerate}

\begin{itemize}
\tightlist
\item
  \textbf{Hint:} group\_by(col1, col2) \%\textgreater\%
  summarise(new\_var = sum(var)) \%\textgreater\%
  arrange(desc(new\_var))
\item
  \textbf{Sugerencia:} Usa \texttt{group\_by(col1,\ col2)}
  \%\textgreater\%
  \texttt{summarise(new\_var\ =\ sum(var))\ \%\textgreater{}\%\ arrange(desc(new\_var))}
\end{itemize}

\begin{enumerate}
\def\labelenumi{\arabic{enumi}.}
\setcounter{enumi}{5}
\tightlist
\item
  \textbf{Identify which Time of day has the highest average opalinus
  count.} Identifica en qué Momento del día se registra el promedio más
  alto de opalinus.
\end{enumerate}

\begin{itemize}
\tightlist
\item
  \textbf{Hint:} Use group\_by(Time) \%\textgreater\%
  summarise(avg\_opalinus = mean(opalinus)) \%\textgreater\%
  arrange(desc(avg\_opalinus))
\item
  \textbf{Sugerencia:} Usa group\_by(Time) \%\textgreater\%
  summarise(avg\_opalinus = mean(opalinus)) \%\textgreater\%
  arrange(desc(avg\_opalinus))
\end{itemize}

\end{document}
